\title{LAC - Assignment 5}

\documentclass[11pt]{article}
\usepackage[utf8]{inputenc}
\usepackage[english]{babel}
%\documentclass[11pt]{article}
\usepackage{setspace} %to set line spacing
\usepackage{fullpage} %uses up more of the page
\usepackage{titlesec} %to edit section titles
\usepackage{graphicx} %to import graphics
%\usepackage{siunitx}
%\usepackage{pdflscape}
%\usepackage{longtable}
\usepackage[title]{appendix}
\usepackage{float}%for the H option in figures
\usepackage{bigstrut} %for excel to latex tabular converter
\usepackage{lscape} %for landscape pages
\usepackage{pdfpages}
\usepackage{amsmath} %\begin{align} - \end{align}
\usepackage[framed,numbered]{mcode}  %imports and formats .m files
\usepackage{multirow}
%\usepackage{tikz}
%\usetikzlibrary{shapes.geometric, arrows}
%\usepackage{iftex}
%\usepackage{array,booktabs}
\usepackage{subcaption}
\usepackage{bm} %Bold in mathmode
%\usepackage{chngcntr}
%\counterwithout{figure}{section}
%\counterwithout{table}{section}
%\counterwithout{figure}{chapter}
%\counterwithout{table}{chapter}
\captionsetup{font=small}
%\usepackage[labelfont=bf]{caption}
%\newcolumntype{C}[1]{>{\centering\arraybackslash}m{#1}}

% BIBLIOGRAPHY SETTINGS
\usepackage{natbib} %for harvard referencing
\bibliographystyle{abbrvnat}
\setcitestyle{numbers,square}

%\cite{berg2009active}  -> [1]
%\citet{berg2009active}  -> Berg et al. [1]
%\citeauthor{berg2009active} -> Berg et al.




\graphicspath{{./Figures/}}

%\myFigure[figure size]{filename.jpg}{Description}{figx}
%\myFigure[]{}{}{}
\newcommand{\myFigure}[4][0.8]{
\begin{figure}[H] 

	\centering
	\setlength\fboxsep{1pt}
	\setlength\fboxrule{0.5pt}

	\includegraphics[width=#1\textwidth]{#2}
	\caption{\small #3}
	\label{#4}
\end{figure}
}

\newcommand{\myFigureSideSide}[8]{
\begin{figure}[H] 

	\centering
    \begin{subfigure}[b]{0.49\textwidth}
    \centering
    \includegraphics[width=1\textwidth]{#1}
    \caption{\small #2}
	\label{#3}
    \end{subfigure}
    \begin{subfigure}[b]{0.49\textwidth}
    \centering
    \includegraphics[width=1\textwidth]{#4}
    \caption{\small #5}
	\label{#6}
    \end{subfigure}   
% 	\setlength\fboxsep{1pt}
% 	\setlength\fboxrule{0.5pt}
	\caption{\small #7}
	\label{#8}
\end{figure}
}



\setlength{\parindent}{0cm} %removes paragraph indentation
%\newcommand{\sr}{\rule[-0.45cm]{0pt}{0.9cm}}


\titleformat{\subsection}{\normalsize\bfseries}{\thesubsection}{1em}{}
\titleformat{\subsubsection}{\normalsize\it\bfseries}{\thesubsubsection}{1em}{}
 \newcommand{\hinfty}{$\mathcal{H}_\infty$~} 
 \newcommand\norm[1]{\left\lVert#1\right\rVert}

\begin{document}

\begin{titlepage}
\includepdf{CoverPage}
\newcommand{\HRule}{\rule{\linewidth}{0.5mm}} % Defines a new command for the horizontal lines, change thickness here

\center % Center everything on the page

%----------------------------------------------------------------------------------------
%	HEADING SECTIONS
%----------------------------------------------------------------------------------------
\begin{figure}[H]
\centering
\includegraphics[width=40mm]{Figures/DTU_Logo.png}
\end{figure}
\vspace{3cm}
% \textsc{\LARGE TU Delft}\\[1.5cm] % Name of your university/college
\textsc{\Large Active Tip Deflection Control For Wind Turbines}\\[0.5cm] 

%----------------------------------------------------------------------------------------
%	TITLE SECTION
%----------------------------------------------------------------------------------------

\HRule \\[0.4cm]
\setstretch{1.3}
{ \huge \bfseries Literature Review}\\[0.4cm] % Title of your document
\setstretch{1}
\HRule \\[1.5cm]


%----------------------------------------------------------------------------------------
%	AUTHOR SECTION
%----------------------------------------------------------------------------------------

\begin{center}
\begin{tabular}{rll}
&\\
\large
Jaime Liew && s141777 \\

\end{tabular}
\end{center}

\vspace{3cm}
~
~

\vspace{2cm}

% If you don't want a supervisor, uncomment the two lines below and remove the section above
%\Large \emph{Supervisor:}\\
%\textcolor{red}{John \textsc{Smith}}\\[3cm] % Your name

%----------------------------------------------------------------------------------------
%	DATE SECTION
%----------------------------------------------------------------------------------------

{\large \today}\\[2cm] % Date, change the \today to a set date if you want to be precise
 
%----------------------------------------------------------------------------------------
%	LOGO SECTION
%----------------------------------------------------------------------------------------

%\includegraphics[width=0.15\textwidth]{figures/DTU.jpg} \\[1cm] % Include a department/university logo - this will require the graphicx package
 
%----------------------------------------------------------------------------------------

\vfill % Fill the rest of the page with whitespace

\end{titlepage}


\section{Introduction}

\section{Literature Review}

\subsection{Background}
Since the advent of modern wind turbine design, the capacity and size of wind turbines has increased dramatically. From the first electricity-producing wind turbine with a rotor diameter of 17m, the design trend has shown huge growth in rotor size with a 180m rotor on the Adwen 8MW platform. As the size of wind turbines continues to increase, new challenges arise in terms of rotor aerodynamics, structural design and control. large rotors require slender, lighter blades to maintain aerodynamic efficiency and to combat the increase in manufacturing costs. This comes at the cost of exacerbated fatigue loads as the blades experiences various forces as they sweep around the rotor plane.
\\~\\
It is well known in the wind industry that loads induced by gravity loads, tower shadow, wind shear, turbulence and gusts can cause significant damage to the blade roots, rotor shaft and the tower \cite{13_Dolan}. By introducing a larger rotor, not only is the turbine more susceptible to this damage, but new instabilities can also be introduced such as aeroelastic flutter, whirling modes, and dynamic coupling such as tower tilt-yaw and blade flap-torsion coupling \cite{7_Berg}.
\\~\\
These problems are typically addressed in the aerodynamic and structural design of the wind turbine to ensure significant damage does not occur by introducing safety factors and increasing blade strength. It is becoming increasingly difficult to both predict the effect of these weaknesses and instabilities as well as effectively eliminate them in the design stage without incurring costs in other aspects of the design, such as increased blade mass. For this reason, the wind industry is focusing more attention on innovative methods for reducing loads. One method, which is the focus of this thesis, is the use of active control methods to actively mitigate fatigue loads in order to extend the lifespan of wind turbine components. In particular, the use of tip deflection sensors in a control system will be investigated.
\\~\\
\subsection{Fatigue Load Mitigation in Wind Turbine Control Systems}
The concept of active load control is not a new concept in the wind industry. Tower and drive train loads can successfully be reduced by extending the basic speed control used in modern variable speed wind turbines.
\subsubsection{Torque Control}
Torque control is used to maintain optimal power output when the wind turbine is operating below rated wind speed. The idea is to balance the aerodynamic and generator torque to achieve a maximum power coefficient $C_p$. A common control strategy to achieve this is to use a proportionality relationship between demanded torque $Q$, and the square of the generator angular velocity $\omega_g$
$$Q = \frac{\pi \rho R^5 C_p}{2\lambda^3 G^3}\omega_g^2$$ \cite{15_bossanyi}

torque control has also been used to introduce side-side tower damping and mitigate drive train torsion vibration. WRITE MORE
\subsubsection{Collective Pitch Control}
Collective pitch control is the basis for turbine power control in modern wind turbines when operating above rated wind speed. Two common objectives are constant torque control or constant power control. Both control objectives are achieved by collectively pitching the turbine blades to control the aerodynamic torque of the rotor. 
\\~\\
It is common to superimpose a high bandwidth controller over the the power/torque controller to decrease the fore-aft tower motion. This is achieved by introducing artificial damping into the tower motion dynamics, reduces fatigue loads in the tower root as well as providing a more stable power output. Control at these frequencies has a much higher bandwidth than the power control loop. Therefore the two controllers are generally superimposed without noticeable interference \cite{15_bossanyi}. \citet{17_Geyler} (check reference), the tower fore-aft damping control is highly coupled with the speed control due to the change in axial wind speed. 

\subsubsection{Individual Pitch Control}
Although collective pitch and torque control methods can reduce oscillations in the tower and drive train, it is ineffective at reducing certain oscillations in the blades. The reason for this is due to azimuth angle dependent loads such as wind shear and tower shadow. Furthermore, each blade experiences different stresses due to variation in the turbulent structure of the wind field. This is the motivation for individual pitch control (IPC). IPC has shown promising results in literature and simulation, and has only recently been introduced into commercial wind turbines, such as the Vestas V164-9.5.
\\~\\
IPC in literature has shown great reductions in flapwise blade loads using a variety of controller designs. \citet{4_trudnowski} achieved an 86\% reduction in flapwise loads using only the rotor angle as an input signal, however the simulation neglected the effects of turbulence. \citet{5_Bossanyi} showed reductions in equivalent fatigue loads in the out-of-plane (OOP) blade root moments, as well as shaft and yaw bearing moments using IPC compared to CPC. \citet{6_Mirzaei} compared model predictive control and PI control for an IPC system, and found comparable reductions in OOP blade root bending moments in stochastic wind speeds based on LIDAR measurements. \citet{14_Selvam} compares IPC systems using PI control as well as LQG control. PI control achieved load reductions at low frequencies, and the LQG controller was able to achieve load reductions at a higher bandwidth, including 2P and 3P, therefore able to reduce loads on non-rotating parts such as the nacelle, yaw bearing and tower. IPC using~\hinfty control was addressed in \citet{1_Lu}, \citet{2_Kanev} and \citet{17_Geyler}, showing not only reductions in OOP blade root bending moments, but also adequate robustness from unmodelled and stochastic behaviour in the system. In all the considered literature, and to the best of the authors knowledge, load reductions were best achieved in OOP blade root moments rather than in plane moments. The reason for the poor performance of edgewise oscillation control is due to the large magnitude of the gravity loading which can not avoided easily \cite{4_trudnowski}.
\\~\\
An issue with IPC is the pitch rate of the blades. Pitch rates required to achieve decent reductions in fatigue load are around $\pm10$ deg s$-^1$, which is considered quite high\cite{15_bossanyi}. However, the required pitching rate decreases with rotor diameter due to the decrease in rotational frequency. This justifies the use of IPC for larger wind turbine models. Higher order harmonic control may not meet the limiting pitch rate requirements, which explains why most papers consider low frequency oscillations (usually up to 3P). \cite{17_Geyler}
\\~\\
There are also mechanical concerns regarding the use of IPC. An increase in wear is expected as the pitch must shift at each rotor rotation. Heat dissipation could also be a concern in the pitch actuators, and should be taken into account in the performance of the actuators at different operating temperatures.


\subsubsection{Active Aerodynamic Load Control}
Active aerodynamic load control (AALC) devices such as trailing edge flaps and micro tabs have also been an active area of research. The advantage of AALC devices is their high bandwidth, allowing for controllability of high frequency dynamics. \citet{7_Berg} and \citet{10_Wilson} have researched the effects of load reduction using trailing edge flaps, showing a 20-32\% reduction in blade root stress, which can allow for a 10\% increase in blade length wihtout exceeding the original equivalent fatigue damage. These papers use tip deflection and tip deflection rate  as controller inputs, making them highly relevant to this the topic at hand. The research in this field is so far restricted to small scale wind turbine models and simulations due to the difficulty and cost of implementing and maintaining such a system on a full scale wind turbine. 
\subsection{Sensing methods}
Real time measurements have become increasingly important for wind turbines and wind farms over recent years, allowing for sophisticated digital control. The stream of data is typically referred to as Supervisory control and data acquisition (SCADA), and can comprise of a variety of novel measurements from components of the turbine. A summary of sensors relevant to active load reduction is described in this section.
\subsubsection{Rotor Azimuth measurement}
Modern wind turbines use rotor angular velocity for the power control loop rather than the rotor azimuth angle. For IPC, a measurement of the rotor azimuth angle is vital in transforming the rotor frame of reference to the fixed tower frame of reference. This is because many load oscillations are dependent on the azimuth angle, such as tower shadow, wind shear and yaw misalignment. 
\\~\\
\citet{4_trudnowski} demonstrated effective control of flap wise loads given only the rotor azimuth angle. Given a precise model of how the wind field varies with rotor angle, for example, \citet{13_Dolan}, this type of control can be possible. However,the wind field can vary dramatically for a given site, and therefore a controller using only rotor azimuth angle would likely be ineffective in the real world. Additionally, fluctuations due to turbulence can not be effectively reduced without additional sensing and controller feedback. Nevertheless, the use of rotor azimuth sensing in conjunction with additional sensors is relevant to this project.

\subsubsection{Wind measurements}
Cup anemometers and wind vane are the standard, IEC-approved method for taking wind speed and direction measurements, however they are unable to capture high frequency components of the wind field as well as spacial variation. As with the power control, load reduction controllers can benefit from the low frequency cup anemometer measurements for gain scheduling. Wind fields can vary greatly over time and space, more advanced sensors may be beneficial to observe how the wind field varies with rotor azimuth angle. LIDAR measurements can provide wind field data at both a high frequency and spacial resolution, and can provide useful information to a control system regarding gusts and field variations before they are experienced by the turbine. \citet{6_Mirzaei} Demonstrated that model predictive control methods were more effective than PI control when used with upstream LIDAR measurements especially in transitioning wind conditions. 
\\~\\
Although the wind measurement techniques is beyond the scope of this project, it should be noted that tip deflection sensors could provide useful information in terms of gust control when used in conjunction with wind speed measurements. \citet{2_Kanev} uses blade root bending moment measurements, both in-plane and out-of-plane, in order to better estimate the blade-effective wind speed, and to additionally detect gust events. An extended Kalman filter observer was used as a nonlinear observer for this study.
\\~\\
...
\subsubsection{Tower Acceleration}
Tower motion, both in the fore-aft and side-side direction, is the source of fatigue loading in the tower. Additionally, the fore-aft dynamics of the tower is coupled with the power output of the turbine due fluctuations in the effective wind speed at the rotor. For this reason, tower damping is often a control objective using pitch and torque control. In order to effectively reduce this oscillation, measurements of the tower acceleration are generally used as feedback to the control system. This can easily be performed with an accelerometer in the nacelle of the wind turbine \cite{15_bossanyi}.
\subsubsection{Strain Gauges}
Commonly used in literature for IPC. No apparent reason for using strain gauges over tip deflection sensors. will write more here.
\subsubsection{Tip Deflection Sensors}
Tip deflection sensors are not commonly cited in literature. Bossanyi briefly mentions the possibility of using accelerometers in the blade tips as an alternative to strain gauges, and also mentions the difficulty of maintaining such sensors due to inaccessibility of the blade tips. \cite{5_Bossanyi} \cite{15_bossanyi}. \citet{7_Berg} and \citet{10_Wilson} use tip deflection sensors in their turbine controller designs, however they focus on active flap control. To the best of the author's knowledge, no research has been performed on IPC using tip deflection sensors instead of strain gauges. using \textcolor{red}{A whitepaper on the cost benefits and basic implementation of the iRotor would be very nice in this section}.
\subsection{Wind Turbine Modelling}
\textcolor{red}{will elaborate on this section}
\subsubsection{Structural}
Something on Euler-Bernoulli beam theory vs Timoshenko, and that HAWC2 uses Timoshenko. Something about the simplified structural models(usually 1 mode for blades and tower) used in literature, and the different assumptions they make and why. Usually a simplified model will suffice and provides transparency to the controller.
\subsubsection{Aerodynamic Modelling}
Aerodynamic models can be complicated (CFD, actuator models), difficult to implement in a control system. Can use advanced models to get linear approximations of system dynamics. Talk about how this is addressed in literature for IPC.
\cite{11_Wang}

\subsection{Observer design}
Kalman filter \cite{15_bossanyi}, extended kalman filter \cite{2_Kanev}. Implementation, and inclusion of colored wind spectrum \cite{14_Selvam}. 
\\
A fundamental aspect of a wind turbine controller is the state observer. The state observer recreates the structural and aerodynamic state of the wind turbine system. A good observer is able to improve controller response by providing a short term prediction how the state will change. It is able to track the state precisely over different working conditions, and in the presence of noise, both from the measurements and from the system dynamics itself. A standard and highly effective observer design is typically the Kalman filter. A Kalman filter is able to DESCRIBE THE KALMAN FILTER. In linear control theory, a Kalman filter is straight forward to design if the process and measurement noise can be known a priori. The Kalman filter can be extended to nonlinear systems by using an extended Kalman filter, which uses a linearized approximation of the current operating point. Something about the unscented Kalman filter. Something about a second order extended Kalman filter (augmented Kalman filter) (maybe not?).
\\~\\
An alternative and novel approach to state observation is using machine learning algorithms. WRITE ABOUT THIS
\subsection{IPC Controller design}
As the control problem at hand is nonlinear, there are a number of controller designs which have been proposed in literature regarding IPC. One common step is performing a transformation to convert the system into a time invariant system. This is called the Coleman transformation, which is a special case of the the multi-blade coordinate transformation, and also known as the d-q transformation borrowed from electrical machine theory. The Coleman transformation removes the dependency of rotor azimuth angle from the system, allowing for linear control methods to be used. More specifically, the Coleman transformation converts the flapwise and edgewise loads into tower tilt and yaw moment, effectively representing 3P loads (or is it 1P?) into 0P loads. In addition to the Coleman transformation, many control ideologies have been adopted in literature. Some of the more common methods are outlined in the following section.

\subsubsection{PI Control}
PI control is the most basic control method used in literature. PI control is a linear control method which provides a control signal proportional to to the error signal and its integral over time. When used in IPC, it is assumed that the yaw and tilt moment decomposition resulting from the Coleman transformation are decoupled. This allows the tilt and yaw to be treated independently with two separate PI controllers. PI control is often used as a benchmark for testing other control methods in IPC. Despite the simplicity of the control method, it shows comparable results to more advanced models at low bandwidths and is more transparent in its implementation \cite{6_Mirzaei}, \cite{14_Selvam}. 
\\~\\
The underlying assumption of using PI control for IPC is the decoupling between the tilt and yaw degrees of freedom. In reality, this is not the case. \citet{1_Lu} presents the weaknesses of the decoupling assumption. Additionally, performing the Coleman transformation and its inverse causes a frequency shift in the system dynamics which should be taken into account in the controller design. 
\subsubsection{Linear-Quadratic-Gaussian Regulator}
A linear-quadratic-Gaussian regulator (LQG) is a multi-input-multi-output (MIMO) control method that provides optimal control to uncertain linear system with white noise disturbances. A LQG is a combination of a linear-quadratic regulator (LQR) with a Kalman filter as a state estimator. For a linear system, the LQR and the Kalman filter can be designed independently of each other. The controller is able to act on multiple sensor inputs and achieve multiple control objectives by changing the weights of a cost function.  \citet{14_Selvam} and \citet{5_Bossanyi} use LQG, and it works pretty well (elaborate more).
\\~\\
The weakness of LQG is that it is a linear controller, which means that it may not perform well on the nonlinear dynamics of a wind turbine. This can usually be addressed by linearising the system about control points and with gain scheduling.
\subsubsection{\hinfty Loop Shaping}
\hinfty loop shaping is a control technique conducted in the frequency domain, where the frequency response of a system can be adjusted through optimisation. \hinfty loop shaping had the advantage of being able to include robustness in its control objective. My understanding of \hinfty loop shaping is a bit rusty. Will elaborate later. 

Used in \cite{1_Lu} to overcome the tilt-yaw coupling problem in PI control. Used in \citet{17_Geyler} on a simple wind turbine model. The unmodelled behaviour is taken into account in the robustness of the \hinfty controller. Also successfully used in \citet{2_Kanev}.
\subsubsection{Model Predictive Control}
Model predictive control used in \cite{6_Mirzaei}.
\subsubsection{Machine Learning}

Neural networks and machine learning. \cite{8_Wang} \cite{3_Treiber}. bad because black box, and does not add any benefits over modern control theory. the motivation is lacking according to \citet{15_bossanyi}

...Cyclic pitch control - Selvam.



% The unique aspect of this paper is in the use of tip deflection sensors instead of a strain gauge at the blade roots. This has rarely been investigated in literature. \citet{7_Berg} addresses the use of tip deflection readings as an input for a turbine control system using active flaps. DESCRIBE RESULTS. Although active flaps can provide a high frequency response for a turbine controller, they have yet to be used in commercial turbine models, and are typically reserved for academic research. Active flaps are expected to increase operating and maintenance costs due to the additional moving parts required. XX mentions the possibility of using tip deflection readings as an alternative input. Apart from these references, tip deflection sensors have not been largely considered in literature to the best of the author's knowledge. It follows that the use of tip deflection sensors in conjunction with IPC has not been explored, and will remain the focus of this paper.
% \\~\\
\textcolor{red}{Not sure where to write a section on the types of loads, or what to cite.
Two types of loads are typically considered in engineering design: fatigue load and ultimate load. Fatigue loads occur as a result of structural oscillations over a long period of time. Calculations relating to fatigue largely consider the magnitude of the loading rather than the frequency. The rainflow counting algorithm developed by XX, is commonplace in engineering in distinguishing the the number of cycles of different magnitudes in a timeseries. Using the rainflow counting results, a short term equivalent load can be determined using BLAH BLAH. An equivalent load is a way of representing the fatigue load, which may consist of various amplitude oscillations at different frequencies, as a single sinusoidal oscillation with a fixed frequency (typically 1Hz). One control objective could be to minimise the short term equivalent load of the wind turbine blades over the operating lifetime.
\\~\\ 
The second type of load that is considered is the ultimate load, which is the largest expected load to occur over a short period of time. Extreme loads which exceed the ultimate strength of a component can lead to instantaneous catastrophic failure of a wind turbine, and can occur due to extreme winds or gusts. IEC standards consider the most extreme load within a 10 minute period which is expected to occur over 50 year time period. Extreme analysis, such as the use of a Gumbel distribution is commonly used to predict the 50 year extreme load based on limited time series data.}
\\~\\
\textcolor{red}{Take into account torsion?
\citet{7_Berg} says flap bending - torsion coupling increases with rotor sizze. ignores torsion as impact is minor. }
\\\textcolor{red}{Is cyclic pitch control the same as IPC? Should I make a distinction?}
\\
\textcolor{red}{Is a sensor to measure azimuth angle the same as the sensor used to measure rotor speed? }\\
\textcolor{red}{What are the benefits of using tip deflection sensors}

\section{Theoretical Analysis}
In this section, the behaviour of the tip deflection is analysed and formulated in a way which is suitable for controller design. First, a model will be presented which describes how the root bending moment of turbine blades behaves. Then, the model will be modified to represent tip deflection.
\subsection{Blade Root Bending Moment Modelling}
A modified blade model used in \cite{1_Lu} and ?? will be used to represent the influences and effects on blade root bending moments on wind turbine blades. Each blade is modelled using the first flapwise bending mode. The wind turbine model in Figure \ref{fig:turb_model} considers the inputs from the blade pitch angle, $\theta_i$, and the relevant output in this case is the flapwise root bending moment of each blade, $M_i$. The collective pitch control is also included in the model diagram as it provides an influence on the blade pitch. All other disturbances on the turbine including turbulence, aerodynamic and structural effects, are encapsulated in the term $D(t)$.

FIGURE GOES HERE fig:turbine_model

Linearizing this model about an operating point with a fixed rotor speed, the collective pitch controller will provide a constant pitch demand for the turbine. 

\subsection{Numerical Assessment}
The transfer function between tip deflection and root bending moment is investigated using numerical analysis methods on simulated time series data. Time series data for tip deflection and bending moment were generated using HAWC2 over a range of windspeeds from 4m/s to 26m/s, and for both normal and extreme turbulence models. 
\\~\\
ASSUMPTIONS (tower shadow, wind shear, etc)
\\~\\
The system identification package in MATLAB is used to estimate the transfer function. A fit percentage is given to each fit using the Normalized Root Mean Squared Error (NRMSE), defined as
$$\text{NRMSE}=\left(\frac{\norm{x_{\text{meas}}-x_{\text{model}}}_2}{\norm{x_{\text{meas}}-\overline{x_{meas}}}_2}\right)$$
Results show that a transfer function with 2 poles and 1 zero is adequate for representing the relationship between tip deflection and bending moment compared to other linear systems. This implies a relationship between the tip deflection rate and a second order system describing root bending moment:
$$a\ddot{M} + b\dot{M} + cM = \dot{x}_{td}$$


\pagebreak

\section{Methodology}
Design objectives:
\begin{itemize}
    \item Load reductions in realistic operating conditions.
    \item Control actions within the limitations of the pitching actuator.
    \item Implementable in discrete time.
\end{itemize}
\subsection{Blade Model}
\subsubsection{Azimuth Based Control}
\subsection{Filter Design Using Zero-Pole Placement}
\subsubsection{Second Order Low Pass Filter}
\subsubsection{Lead Compensator}
A lead compensator is a component of a control system which compensates for undesired lag in a system. A lead compensator is an alternative method to PID control methodology in achieving a desired frequency response of a controller. The transfer function of a lead compensator consists of a pole-zero pair on the real axis of the s-plane:

$$C(s) = K\frac{s-z}{s-p}$$
where $z$, $p$ are the zero and pole of the compensator. For a lead compensator, $|z| < |p|$. The amount of phase lead and the frequency at which the compensation occurs can be specified by the choice of $z$ and $p$. It is convenient to reformulate the compensator transfer function in terms of a time constant, $T$, and attenuation constant, $a$, by substituting $p = 1/T$, $z = 1/aT$, $K=a$, which leads to:
$$C(s) = \frac{1+aTs}{1+Ts}$$
The maximum amount of phase lead, $\phi_m$, occurs at a frequency, $\omega_m$, and are related to $T$ and $a$ by the relations:

$$\sin\phi_m = \frac{a-1}{a+1}$$
$$\omega_m = \frac{1}{T\sqrt{a}}$$
The value of using a lead compensator in the tip deflection controller is that an arbitrary phase lead angle can be set at a target frequency to help overcome the large amount of lag in the system. Given a desired phase lead at a given frequency, $T$ and $a$ can be found, from which the continuous transfer function in terms of $z$, $p$ and $K$ can be found. 
\subsection{Bandpass Filter ?}

\subsection{Controller Discretisation}
Although control design is performed in continuous time, the controller for this application is implemented digitally. For this reason, the continuous controller must be transformed to be implemented as a discrete time controller. The advantage of designing an LTI controller is that there is a large body of knowledge for performing continuous to discrete transformations. 
\\~\\
It is not possible to produce a discrete system which perfectly matches the frequency and time domain performance of its continuous counterpart. For this reason, many methods exist. The Impulse invariant method, zero order hold method, and first order hold method are examples if continuous-to-discrete methods which produce discrete systems which exactly match the time-domain response if a continuous system, namely the impulse, step, and ramp response respectively. The zero order hold method, for example, is the default method used in MATLAB's \texttt{c2d} function. Despite its common use, time-domain invariant methods do not preserve the frequency domain response of a system. As the frequency response of an IPC controller is of high importance, a different discretisation method is sought after.
\\~\\
Two methods were investigated which better preserve the frequency response of the system: pole-zero mapping and the bilinear transform. Both methods showed comparable results due to the high sampling frequency of the controller compared to the frequencies of the wind turbine system.
\subsubsection{Zero-pole mapping}
Pole-zero mapping involves directly mapping the continuous poles and zeros, $s_i$ from the s-domain to equivalent locations in the z-domain, $z_i$, using the transformation $z_i = e^{s_iT_s}$ where $T_s$ is the sampling time of the discrete system. Expressed another way, a continuous transfer function with poles and zeros, $p_i$ and $z_i$:
$$H_a(s) = K_a\frac{\prod_{i=0}^{N}(s-z_i)}{\prod_{i=0}^{N}(s-p_i)}$$
is transformed to a discrete transfer function:
$$H_d(z) = K_d\frac{\prod_{i=0}^{N}(1-e^{z_iT_s}z^{-1})}{\prod_{i=0}^{N}(1-e^{z_iT_s}z^{-1})}$$
TODO
\subsubsection{Bilinear Transformation}

\subsubsection{1P Control Design}
The first feedback controller to successfuly reduce fatigue loads in this project attenuates disturbances at 1P only. This controller consists of a second order low pass filter with a cutoff frequency set at 1P, as well as two identical lead compensators:

$$C_{1P}(s) = K\underbrace{\frac{1}{s^2 + 2\zeta\omega s + \omega^2}}_\text{Low pass filter}\underbrace{\frac{(1-aTs)^2}{(1-Ts)^2}}_\text{Lead compensator}$$

The parameters of this controller are defined in Table XX. Figure XX shows the bode plot of the controller. Note that the controller magnitude response peaks at 1P and is highly attenuated at all other frequencies. Furthermore, the phase response at 1P has a phase lead of XXdegrees to account for the lag in the system. The full continuous transfer function of this system is:
$$continoustransfer function$$
Using the bilinear transform, the controller is discretised to XX:


$$discrete transfer function$$
\begin{table}[]
\centering
\caption{My caption}
\label{my-label}
\begin{tabular}{r|l}
\textbf{Parameter}           & \textbf{Value} \\ \hline
Damping coefficient, $\zeta$ & 0.05           \\
Cutoff frequency, $\omega$   & 1.005          \\
Time constant, $T$           & 0.184          \\
Attenuation constant, $a$    & 29.11         
\end{tabular}
\end{table}
\subsection{HAWC2 integration}
Implementation of a digital IIR filter for each blade.
pseudocode of controller.

\section{Results}
\subsection{Quantifying fatigue loads}
The turbine simulations produce complicated time histories of the loads experienced in the various components of the turbine. In order to quantify and compare the fatigue damage experienced in these components, a load spectrum is calculated. A load spectrum decomposes a complicated stress history into stress cycles of varying amplitude. This is achieved using rainflow counting\textbf{source to matsuishi and endo}.The fatigue loads in the turbine components can be quantified by calculating the equivalent load at 1Hz. The equivalent load is the amplitude of a 1Hz oscillating load which produces the same amount of fatigue damage to a component as a mixed load spectrum. It is a way of comparing different load spectra of the same component. The Equivalent load, $S_{eq}$, is calculated for $N_{eq}$ cycles as follows:

$$S_{eq} = \left(\frac{\sum_i S_iN_i^m}{N_{eq}}\right)^{\frac{1}{m}}$$

Where $S_i$ is the $ith$ load cycle amplitude, $N_i$ is number of full cycles at $S_i$, $m$ is the material Wh\"{o}ler curve exponent of the component in question. In this analysis, a Wh\"{o}ler curve exponent of 4 and 10 is used for steel (tower, rotor shaft, etc) and composite materials (ie. rotor blades) respectively. $N_eq$ is the number of cycles experienced of load $S_{eq}$. For 600 second simulations, which is the case for this analysis, A 1Hz equivalent load requires $N_{eq}=600$. 
\pagebreak
\bibliography{0_Bibliography}

\end{document}
